\section{相关工作}
\label{sec:related}

已有工作通过无导数策略学习或在线规划来学习视觉控制中的潜在动力学，也有工作用多步预测增强无模型智能体，或利用 Q 值和多步奖励的解析梯度；这些方法通常面向低维任务。相比之下，Dreamer 纯粹依靠潜在想象，并通过解析梯度高效学习视觉控制中的长时程行为。

\paragraph{基于潜在动力学的控制}

E2C~\citep{watter2015e2c} 和 RCE~\citep{banijamali2017rce} 将图像嵌入到紧凑空间中，并在该空间内向前预测，用于求解简单任务。World Models~\citep{ha2018worldmodels} 以两阶段方式学习潜在动力学，并在想象中演化线性控制器。PlaNet~\citep{hafner2018planet} 联合学习这些组件，并通过潜在空间中的在线规划解决视觉运动任务。SOLAR~\citep{zhang2018solar} 通过潜在空间中的引导策略搜索解决机器人任务。I2A~\citep{weber2017i2a} 将想象轨迹交给无模型策略，而 \citet{lee2019slac} 和 \citet{gregor2019beliefstate} 则学习信念表征，以加速无模型智能体。

\paragraph{想象的多步回报}

VPN~\citep{oh2017vpn}、MVE~\citep{feinberg2018mve} 和 STEVE~\citep{buckman2018steve} 从重放缓冲区中学习动力学，用于多步 Q 学习。AlphaGo~\citep{silver2017alphago} 将动作和状态价值预测与规划结合起来，但假设可以访问真实动力学。同样假设可访问动力学的 POLO~\citep{lowrey2018polo} 通过学习价值集成来进行探索。MuZero~\citep{schrittwieser2019muzero} 学习特定任务的奖励和价值模型以解决困难任务，但需要大量经验。PETS~\citep{chua2018pets}、VisualMPC~\citep{ebert2017visualmpc} 和 PlaNet~\citep{hafner2018planet} 使用无导数优化进行在线规划。POPLIN~\citep{wang2019poplin} 通过自我模仿改进在线规划。\citet{piergiovanni2018dreaming} 使用潜在动力学模型，通过想象学习机器人策略。神经网络梯度规划已在小规模问题上得到验证\citep{schmidhuber1990diffmodel,henaff2018planbybackprop}，但扩展到大规模任务仍有困难\citep{parmas2019pipps}。

\paragraph{解析价值梯度}

DPG~\citep{silver2014dpg}、DDPG~\citep{lillicrap2015ddpg} 和 SAC~\citep{haarnoja2018sac} 利用学得的即时动作价值梯度，通过经验回放学习策略。SVG~\citep{heess2015svg} 使用单步模型预测的解析价值梯度，降低无模型 on-policy 算法的方差。\citet{byravan2019ivg} 的同期工作在导航和操纵任务中使用潜在想象与确定性模型。ME-TRPO~\citep{kurutach2018modeltrpo} 通过对本体感知输入的预测回报求梯度，加速原本无模型的智能体。DistGBP~\citep{henaff2017planbybackprop,henaff2019planbybackprop} 在简单任务中使用模型梯度进行在线规划。
